% ===== 原文 pp.32–33（PDF pp.33–34）=====

% [原文 p.32]
\begin{figure}[htbp]
\centering
\includegraphics[width=0.9\textwidth]{fig5.png}
\figtitle{5}{$\mathrm{H}_2$（左）与 $\mathrm{H}_3$（右）的 RG 流。}
\end{figure}

其中有些流会穿过相 I。我们来讨论能否把 $\mathbb{F}_n$ sigma 模型认作模型 $\mathrm{H}_n$ 的一部分。我们写出 FI-theta 参数为 $q_1=(\Lambda/\mu)^{2-n}$、$q_2=q(\Lambda/\mu)^2$，其中 $q$ 与 $\Lambda$ 是重整化群不变量。当 $q_1\ll 1$ 且 $q_2\ll 1$ 时理论处于相 I，也就是处在能标窗口 $q^{1/2}\Lambda\ll\mu\ll\Lambda$ 内，而这只有在 $q\ll 1$ 时才可能。如果我们希望保留 $\mathbb{F}_n$ 的两个 Kähler 参数，就应该让 $q$ 与 $\Lambda$ 都保持有限；但若放弃其中一个或两个，则可以取某些极限。例如，取极限 $q\to 0$ 并保持 $\Lambda$ 固定（同时考察 $|\phi_3|^2\sim\zeta^2$ 的场空间）会使 $\phi_3$ 这一对场与第二个 $U(1)$ 规范群退耦，把模型约化为 $\mathrm{T}^{U(1)}_{N,d}(0)$（取 $N=2$、$d=n$）——即它聚焦到轨形点及其消解 $\mathrm{tot}[\cO_{\mathbb{P}^1}(-n)]\to\C^2/\Z_n$。真空方程可以用 $z:=\sigma_2-n\sigma_1$ 写成
\begin{equation}
    (z^2-q\Lambda^2)^2=n^2\Lambda^{2-n}z^{n+2};\qquad
    \sigma_1=\frac{q\Lambda^2-z^2}{nz},\qquad
    \sigma_2=\frac{q\Lambda^2}{z}. \tag{4.38}
\end{equation}
% [Tutor Note: (4.38) 可逐步验证：由 q_1=(\Lambda/\mu)^{2-n} 得 q_1\mu^{2-n}=\Lambda^{2-n}，代入 (4.35)
% 第一式即 \sigma_1^2=\Lambda^{2-n}z^n；(4.35) 第二式即 \sigma_2 z=q\Lambda^2，正是 (4.38) 中
% \sigma_2=q\Lambda^2/z 的来源。再将 (4.38) 的 \sigma_1=(q\Lambda^2-z^2)/(nz) 平方并与 \Lambda^{2-n}z^n 比较，
% 即得第一式 (z^2-q\Lambda^2)^2=n^2\Lambda^{2-n}z^{n+2}。n=2 时右端为 4z^4，方程关于 z 是四次式，故有
% 四个解，与下文一致。]
当 $q\ll 1$ 时，有四个解满足 $z^2\sim q\Lambda^2$，对应 $(\sigma_1,\sigma_2)\sim(\epsilon_1 q^{\frac n4}\Lambda,\epsilon_2 q^{\frac12}\Lambda)$（其中 $\epsilon_1,\epsilon_2\in\{\pm1\}$）；还有 $(n-2)$ 个解满足 $1\sim n^2\Lambda^{2-n}z^{n-2}$，对应 $(\sigma_1,\sigma_2)\sim(-z/n,q\Lambda^2/z)$。在 $q\to 0$ 且 $\Lambda$ 固定的极限下，前四个解趋于 $(\sigma_1,\sigma_2)\to(0,0)$，而后 $(n-2)$ 个解趋于 $(\sigma_1,\sigma_2)\to(\tilde z,0)$，其中 $\tilde z^{n-2}=(-n)^{-n}\Lambda^{n-2}$。在模型 $\mathrm{T}^{U(1)}_{2,n}(0)$（$n>2$）中，$\sigma=0$ 处的真空对应于消解 $\mathrm{tot}[\cO_{\mathbb{P}^1}(-n)]$，而位于 $\sigma^{n-2}=(-n)^{-n}\Lambda^{n-2}$ 的 $(n-2)$ 个真空是有质量的 Coulomb 真空。这表明前四个解对应于 $\mathbb{F}_n$ sigma 模型的真空。然而，如果我们希望保留 $\mathbb{F}_n$ 的两个 Kähler 参数，就没有任何极限能只挑出这些真空而把其余真空送到无穷远。因此，我们可以辨认出对应于 $\mathbb{F}_n$ sigma 模型的真空，却无法把剩余部分退耦。也就是说，我们无法从 $\mathrm{H}_n$ 中分离出“纯粹的 $\mathbb{F}_n$”。
% [Tutor Note: 本页（p.32）原文有几处拼写错误：“resuolution”应为“resolution”，“remaing”应为
% “remaining”，“there is no limit that select”中的“select”应为“selects”；译文按正确含义翻译，不另改动。
% 另可用量纲核对 (4.38)：z 与 \Lambda 同量纲（z^2\sim q\Lambda^2，而 q=\ee^{-t} 无量纲），左端
% (z^2-q\Lambda^2)^2 量纲为（质量）^4，右端 n^2\Lambda^{2-n}z^{n+2} 为（质量）^{2-n+n+2}=（质量）^4，自洽。]

% [原文 p.33]
\subsection{矩多面体与环面簇}
我们来考虑具有连通阿贝尔规范群、且没有超势的一般模型：
\begin{align*}
    G &= U(1)^k,\\
    V &= \bigoplus_{i=1}^{N}\C(Q_i)\ni\phi=(\phi_1,\ldots,\phi_N),\\
    W &= 0,\\
    \tW &= -\sum_{a=1}^{k}t^a\sigma_a,
\end{align*}
其中 $Q_i=(Q_i^1,\ldots,Q_i^k)\in\Z^{\oplus k}$ 是 $G=U(1)^k$ 的一个权。我们假定 $G$ 在 $V$ 上的作用是忠实的。$\CP^{N-1}$ 模型、模型 $\mathrm{T}^{U(1)}_{N,d}(0)$ 与 $\mathrm{H}_n$ 都是它的特例。由忠实性，$G$ 可以看作 $U(1)^N$ 的一个子群，后者以对角权作用在 $V$ 上。下述群将扮演重要角色：
\begin{equation}
    T_F:=U(1)^N/G. \tag{4.39}
\end{equation}
它是理论的一个整体对称群，我们称之为（环面）味群。作为 $N-k$ 维紧连通阿贝尔李群，它与 $U(1)^{N-k}$ 同构。

假定存在一个几何相或轨形相，在其中任何经典超对称真空处规范群都被破缺到一个有限子群。此时理论在低能下约化为靶空间为真空流形的 sigma 模型，该流形为
\begin{equation}
    X_\zeta=\mu^{-1}(\zeta)/G. \tag{4.40}
\end{equation}
味群 $T_F$ 作为其 Kähler 结构的对称群作用在 $X_\zeta$ 上，特别是作为其辛结构的对称群。事实上，存在一个矩映射 $\mu_F:X_\zeta\to\ii\ft_F^*$。让我们选取一个同构 $f:U(1)^{N-k}\to T_F$：
\begin{equation}
    f:(c_1,\ldots,c_{N-k})\in U(1)^{N-k}\longmapsto
    \Big[\Big(\prod_\alpha c_\alpha^{f_1^\alpha},\ldots,\prod_\alpha c_\alpha^{f_N^\alpha}\Big)\Big]\in T_F. \tag{4.41}
\end{equation}
于是 $\mu_F(\phi)\in\ii\ft_F^*$ 的分量由下式给出：
\begin{equation}
    \mu_F^\alpha(\phi)=\sum_{i=1}^{N}f_i^\alpha|\phi_i|^2,\qquad
    \alpha=1,\ldots,N-k. \tag{4.42}
\end{equation}
% [Tutor Note: (4.41)(4.42) 的下标约定务必看清：f_i^\alpha 中下标 i 是场指标、上标 \alpha 是味群坐标
% （故 \mu_F^\alpha 是第 \alpha 个 U(1) 分量，场 \phi_i 在该 U(1) 下的权正是 f_i^\alpha）。这一点可用
% p.36 的 (4.53) 反查：(4.53) 第一式 \sum_i v_\beta^i f_i^\alpha=\delta_\beta^\alpha 恰是 f^{-1}\circ f=\mathrm{id}；
% 第二式则给出 f\circ f^{-1} 与恒等映射相差规范方向 \sum_a Q_j^a\xi_a^i 的项。若把 (4.41) 的指数误写成
% f_\alpha^i，则 (4.53) 无法成立。]
